This paper has three parts: a study of when a sampler should commit tokens, evidence that revising accepted tokens matters in uniform diffusion, and a rolling sampler that needs no training. Each part has close prior or concurrent work.

\subsection{Rolling Windows for Language}
Several methods move the denoising window forward as tokens finish.
DSB \citep{luo2026dsb} does this for masked diffusion.
Set Block Decoding \citep{gat2025set} fine-tunes an autoregressive model to decode masked tokens ahead with the entropy bound of \citet{ben2026accelerated}.
Concurrent with our work, Clock Diffusion \citep{schiff2026clock} trains a Gaussian diffusion model whose window rolls left to right and commits finished positions to the KV cache.
Its noise schedule depends on position, and our position-based temperature schedule plays the same role at inference time.
\citet{wu2023ar} paired diffusion with the left-to-right bias of language earlier still.
RES uses the same rolling mechanism, and we do not claim the mechanism as new.

\subsection{Freezing versus Revising}
In masked diffusion, an unmasked token stays fixed.
A rolling masked sampler such as DSB therefore never revises a token once it accepts it.
\citet{kim2025fine} fine-tune masked diffusion models so that they can revise tokens they have already revealed.

\subsection{What This Paper Adds}
Freezing accepted tokens in DiffusionGemma lowers macro accuracy from 73\% to 20\% at loose entropy bounds (Section~\ref{sec:too-early}).
A rolling sampler must therefore balance two goals: moving tokens into the KV cache early, and keeping them open to revision.
From the order in which tokens are accepted (Section~\ref{sec:too-late}), we design a commit rule with two properties.
Its minimum shift moves the sampler between block-level commitment and fully incremental commitment, which commits whenever it can.
Its entropy bound still trades accuracy for speed, as in the baseline.
Unlike Clock Diffusion, RES needs no training and runs on an existing model.

\subsection{Rolling Diffusion Beyond Text}
Beyond text generation, the idea of rolling diffusion has been studied extensively in other domains.
Within video generation, rolling diffusion was a natural mechanism to address the temporal nature of video generation and avoid performing diffusion on every frame simultaneously \citep{ruhe2024rolling}.
Many other works include similar mechanisms \citep{chen2024diffusion, song2025history, xie2025progressive}, and \citet{papa2026asyncpatch} carry the idea over to images by assigning different patches different levels of corruption.
